% ═══════════════════════════════════════════════════════════════════════════════
%  CAPÍTULO — Chain of Responsibility
% ═══════════════════════════════════════════════════════════════════════════════

\chapter{Chain Of Responsibility}

\patroninfo{Comportamiento}{223}

% ─── Situación Problema ──────────────────────────────────────────────────────
\section{Situación Problema}

Cuando un anfitrión publica un nuevo alojamiento en Airbnb, la plataforma
debe verificar que el anuncio cumple con una serie de requisitos antes
de hacerlo visible al público. Estos requisitos son heterogéneos: algunos
son puramente técnicos ---las fotografías deben superar un umbral mínimo
de resolución y cantidad---, otros son de contenido ---la descripción
no puede contener información de contacto directo que eluda la
intermediación de la plataforma---, otros son de coherencia de negocio
---el precio base no puede ser anómalamente inferior al promedio de la
zona, lo que podría indicar un error del anfitrión--- y otros son de
cumplimiento legal ---ciertas ciudades exigen un número de licencia
turística válido antes de permitir la publicación---.

Ejecutar todas estas verificaciones en un único componente mezcla
responsabilidades de naturaleza completamente distinta: lógica de
procesamiento de imágenes, análisis de texto, comparación de precios
de mercado y consulta de registros regulatorios conviven en el mismo
lugar. Cualquier cambio en la normativa de una ciudad concreta, o la
incorporación de un nuevo criterio de calidad, obliga a intervenir ese
componente central con el riesgo de afectar verificaciones que no tenían
relación con el cambio.

% ─── Solución ────────────────────────────────────────────────────────────────
\section{Solución}

El patrón Chain of Responsibility organiza cada verificación como un
manejador independiente: \texttt{VerificadorFotos},
\texttt{VerificadorContenidoDescripcion}, \texttt{VerificadorCoherenciaPrecio}
y \texttt{VerificadorLicenciaTuristica}. Cada manejador tiene una
responsabilidad acotada: evalúa el aspecto del anuncio que le corresponde
y, si detecta un incumplimiento, detiene la cadena y devuelve al
anfitrión un motivo de rechazo claro y específico. Si el anuncio supera
su verificación, lo pasa al siguiente manejador sin intervención adicional.

Solo cuando el anuncio recorre la cadena completa sin ser detenido se
procede a su publicación. Cada verificador es completamente ajeno a
los criterios de los demás, por lo que modificar el umbral de calidad
fotográfica no tiene ningún efecto sobre la verificación de licencias.
Incorporar una nueva exigencia ---por ejemplo, un detector de precios
discriminatorios--- se reduce a insertar un nuevo manejador en la cadena
en el punto que corresponda, sin tocar ninguno de los existentes.

% ─── Diagrama de clases ─────────────────────────────────────────────────────
\begin{figure}[H]
    \centering
    % \includegraphics[width=\textwidth]{imagenes/chain_responsibility_diagrama.png}
    \caption{Diagrama de clases del patrón Chain of Responsibility}
\end{figure}


% ─── Roles de las Clases ────────────────────────────────────────────────────
\section{Roles de las Clases}

% Explicar la responsabilidad de cada clase/interfaz

\begin{itemize}
    \item \textbf{Handler:}
    \item \textbf{ConcreteHandler:}
    \item \textbf{Client:}
\end{itemize}


% ─── Main ───────────────────────────────────────────────────────────────────
\section{Main}

% Explicar qué realiza el programa principal

\begin{lstlisting}[language=Java, caption=Main]
 //Codigo del main
\end{lstlisting}


% ─── Salida del Main ────────────────────────────────────────────────────────
\section{Salida del Main}

\begin{lstlisting}[caption=Salida]
 // Salida esperada
\end{lstlisting}


% ─── Clases ─────────────────────────────────────────────────────────────────
\section{Clases}

% ─── Clase 1 ────────────────────────────────────────────────────────────────
\subsection{NombreClase}

\begin{lstlisting}[language=Java, caption=NombreClase]
 
\end{lstlisting}


% ─── Clase 2 ────────────────────────────────────────────────────────────────
\subsection{NombreClase}

\begin{lstlisting}[language=Java, caption=NombreClase]
 
\end{lstlisting}


% ─── Clase 3 ────────────────────────────────────────────────────────────────
\subsection{NombreClase}

\begin{lstlisting}[language=Java, caption=NombreClase]
 
\end{lstlisting}


% ─── Conclusiones ───────────────────────────────────────────────────────────
\section{Conclusiones}

% Explicar:
% - Ventajas del patrón
% - Cuándo usarlo
% - Beneficios obtenidos
% - Posibles desventajas
% - Relación con principios SOLID